% ══════════════════════════════════════════════════════════════════════════════
\section{Discussion}
\label{sec:discussion}
% ══════════════════════════════════════════════════════════════════════════════

\subsection{The Opportunity-Cost Channel in Non-Sahel Africa}

The main finding of this paper is that ML-predicted maize yields have a negative and statistically significant effect on subsequent conflict in the preferred TWFE OLS log-log specification: a 10\% increase in predicted yield is associated with a 6.7\% reduction in log conflict events in the 3-month post-harvest window ($\hat{\beta} = -0.666$, $p < 0.001$). This effect is robust to extending the conflict window to 6 months ($\hat{\beta} = -0.949$) and 12 months ($\hat{\beta} = -1.290$), and survives the addition of a spatial lag control ($\hat{\beta} = -0.377$, $p < 0.001$).

The direction of the result is consistent with the opportunity-cost mechanism advanced by \citet{DalBoDalBo2011}: higher agricultural income raises the return to peaceful economic activity and reduces the relative attractiveness of participation in armed conflict. This interpretation is supported by the conflict-type breakdown, which shows the negative yield effect concentrated in violence against civilians ($\hat{\beta} = -0.524$, $p < 0.001$) and explosions/remote violence ($\hat{\beta} = -0.155$, $p = 0.020$)---categories that reflect deliberate armed-actor decisions rather than spontaneous collective mobilisation.

To contextualise the magnitude: at the sample mean of approximately 3.6 conflict events per district-year, a 10\% yield increase corresponds to a reduction of roughly 0.24 events in the 3-month window---modest in absolute terms, but non-trivial at scale across 1,338 districts over 15 years, and comparable in relative magnitude to effect sizes reported in related studies using rainfall or commodity price instruments \citep{MiguelSatyanathSergenti2004, CoutteynierSoubeyran2014}. The longer-horizon estimates suggest that the cumulative pacifying effect over a full year ($\hat{\beta}_{12\text{mo}} = -1.290$) is approximately twice the 3-month effect, implying that the mechanism operates over the full post-harvest economic cycle rather than at a single point in time.

\subsection{The Resource-Predation Channel in the Sahel}

The most novel and policy-relevant finding of this paper is the systematic reversal of the yield--conflict relationship in the Sahel zone. The interaction model in Table~\ref{tab:aez} yields a differential Sahel effect of $+3.513$ ($p < 0.001$), producing a net Sahel elasticity of $+2.383$: within Burkina Faso, Mali, and Niger, higher predicted yields are associated with \textit{more} conflict rather than less.

This positive relationship is consistent with the resource-predation channel formalised in \citet{Grossman1991} and \citet{DalBoDalBo2011}, and illustrated empirically in the context of extractive resources by \citet{Berman_etal2017} and for commodity price shocks by \citet{DubeVargas2013}. When armed actors can appropriate agricultural surplus at low cost---through market taxation, grain seizures, or control of productive territory---a better harvest raises the expected payoff to predatory activity rather than reducing it. The key condition is the appropriability of the surplus: in the Sahel, where state capacity is limited, borders are porous, and armed groups have established coercive presence in agricultural communities, this condition is plausibly met.

The contrast with the opportunity-cost result is instructive. Both channels can operate simultaneously: the question is which dominates. In Non-Sahel countries---Ethiopia, Malawi, Nigeria, Tanzania---the opportunity cost of agricultural income appears to dominate: higher yields reduce conflict by raising the civilian return to farming. In Sahel countries, the predation channel dominates: higher yields increase conflict by raising the armed-group return to extraction. The interaction term captures this asymmetry precisely, and its magnitude and significance ($p < 0.001$) indicate that the difference is not explained by sampling variation.

This finding complements the evidence in \citet{DubeVargas2013}, who show that positive commodity price shocks to oil---an easily appropriable resource---increase conflict in Colombia, while positive shocks to labour-intensive crops reduce it. Agricultural surplus in the Sahel occupies an intermediate position: it is partially appropriable (through grain market control and direct seizure), making the sign of the yield--conflict relationship depend on the institutional and security environment. The results here provide cross-national evidence consistent with this mechanism, extending it to the African agro-pastoral context.

\subsection{Conflict Type and Mechanism}

The disaggregation by conflict type in Table~\ref{tab:conflict_types} provides additional evidence on the mechanism. The null result for riots ($\hat{\beta} = -0.005$, $p = 0.887$) is informative in two respects. First, it rules out a broad economic discontent interpretation: if higher yields reduced conflict through a generalised income effect that pacified populations, one would expect a negative effect on riots as well as on armed violence. The absence of any riot effect suggests that the mechanism operates specifically at the armed-actor margin rather than through civilian grievance or collective action.

Second, the null riot result is consistent with the observation that riots in Sub-Saharan Africa are predominantly urban phenomena driven by urban food prices, political contention, and ethnic mobilisation \citep{BlattmanMiguel2010}. Rural yield shocks affect farm-gate prices and rural household incomes, but their transmission to urban food prices and urban political dynamics is attenuated by marketing margins, storage costs, and the structure of agricultural value chains. The absence of a riot effect therefore also reflects the limited spatial transmission of the rural income channel to urban conflict.

The significant negative effect on violence against civilians ($\hat{\beta} = -0.524$, $p < 0.001$) is the largest point estimate in the conflict-type analysis and deserves emphasis. Violence against civilians is the category most directly associated with predatory behaviour by armed groups---it includes deliberate attacks, abductions, looting, and forced displacement that are instrumentally motivated by resource extraction or coercive control. A reduction in this category when yields are higher, in a specification that controls for population and district fixed effects, suggests that the profitability of civilian targeting declines when agricultural productivity is high. This could reflect either that civilian populations are better able to resist or negotiate in productive periods, or that armed groups substitute toward other activities when the expected resistance to predation increases.

\subsection{Limitations}

Several limitations should inform the interpretation of these results.

\textit{Yield model performance.} The XGBoost model achieves a spatial $R^2$ of 0.35, which is a substantial improvement over single-sensor approaches but still implies that 65\% of the cross-sectional variation in observed yields is unexplained. Year-to-year within-district variation in predicted yields is further constrained by the relatively smooth temporal dynamics of satellite-derived vegetation composites. If the true yield signal is measured with error, coefficient estimates will be attenuated toward zero, and the true elasticity may be larger in absolute terms than what is reported here. Measurement error also affects the conflict-type and AEZ heterogeneity analysis: if prediction error is correlated with country or zone, estimates of the interaction effect may be biased.

\textit{Crop coverage.} The yield model targets maize as the primary food crop. Maize is the dominant staple in most of the study countries, but sorghum, millet, and cassava are important in parts of the Sahel and in Nigeria's northern zones. The predicted yield therefore captures the productivity of one component of the agricultural economy, and the estimated elasticity may not generalise to households or districts where other crops dominate.

\textit{Causal interpretation.} The use of predicted rather than observed yields is intended to mitigate reverse causality: conflict affects actual output, but the satellite vegetation signal and weather realisations that drive predictions should be less directly affected by conflict at the district level. This argument is plausible but not conclusive. If large-scale conflict causes agricultural abandonment at spatial scales visible to satellite sensors---for example, through mass displacement that empties entire districts---then the predicted yield may itself be partially endogenous to conflict. This concern is most relevant for Ethiopia during the 2020--2022 Tigray conflict and for parts of Mali and Burkina Faso during periods of sustained displacement. The per-country and AEZ results should be interpreted with this caveat in mind.

\textit{Sahel classification.} The Sahel/Non-Sahel classification is a country-level proxy that does not capture within-country agro-ecological heterogeneity. Northern Nigeria, for example, shares many characteristics with the Sahel zone, while southern Mali has higher rainfall than the national average. A more granular AEZ classification using gridded data---for example, from the FAO's Global Agro-Ecological Zones (GAEZ) database---could provide finer-grained tests of the predation mechanism. The country-level interaction is a first-order approximation and should be treated as indicative rather than definitive.
